% =====================================================================
%  Guía práctica: Domina las tablas de contenido avanzadas en ConTeXt
%  Mismo sistema de diseño que las guías anteriores.
%  Compilar con:  context tocavanzado.tex
% =====================================================================
\enabledirectives[overloadmode=error]
%\environment env_tutorial
% ---------------------------------------------------------------------
% 1. Tipografía y colores
% ---------------------------------------------------------------------
%\definefontfamily [mainface] [rm] [Pagella]
%\definefontfamily [mainface] [ss] [Iwona]
%\definefontfamily [mainface] [tt] [ModernMono]
%\setupbodyfont     [mainface,11pt]
%\mainlanguage[es]
%
%\definecolor [inkblue]    [r=0.11, g=0.16, b=0.32]
%\definecolor [accent]     [r=0.72, g=0.20, b=0.20]
%\definecolor [accentb]    [r=0.10, g=0.35, b=0.45]
%\definecolor [resultgreen][r=0.13, g=0.42, b=0.24]
%\definecolor [softgray]   [r=0.45, g=0.45, b=0.45]
%\definecolor [panelbg]    [r=0.96, g=0.96, b=0.94]
%\definecolor [resultbg]   [r=0.94, g=0.97, b=0.94]
%\definecolor [codebg]     [r=0.965,g=0.965,b=0.98]
%\definecolor [codeframe]  [r=0.80, g=0.80, b=0.85]
%
%\setupcolors [state=start]
%
%% ---------------------------------------------------------------------
%% 2. Página y márgenes
%% ---------------------------------------------------------------------
%\setuppapersize   [letter]
%\setuplayout      [
%  topspace=2cm, backspace=3cm,
%  header=1cm, footer=1cm,
%  width=middle, height=middle,
%  leftmargin=0cm, rightmargin=0cm,
%]
%\setupwhitespace [medium]
%\setupindenting  [never]
%
%% ---------------------------------------------------------------------
%% 3. Encabezados de sección
%% ---------------------------------------------------------------------
%\definehead [tutorialsection] [section]
%
%\setuphead [tutorialsection]
%  [style=\bf,
%   color=inkblue,
%   number=yes,
%   numbercolor=accent,
%   numberstyle=\bf]
%
%\setuphead [subsection]
%  [style=\bf, color=accentb]
%
%\setuphead [subsubsection]
%  [style=\bf, color=accentb]
%
%% ---------------------------------------------------------------------
%% 4. Bloques de código
%% ---------------------------------------------------------------------
%\definetyping  [contextcode] [
%  option=TEX,
%  space=on,
%  style=\tt,
%]
%
%\setupbackgrounds [contextcode] [background=color, backgroundcolor=codebg]
%
%\setuptyping [contextcode] [
%  frame=on,
%  framecolor=codeframe,
%  rulethickness=0.4pt,
%  leftmargindistance=1em,
%  rightmargindistance=1em,
%  topspace=0.6em,
%  bottomspace=0.6em,
%]
%
%% ---------------------------------------------------------------------
%% 5. Cajas: Problema, Solución, Explicación, Observación, Resultado
%% ---------------------------------------------------------------------
%\defineframedtext
%  [problema]
%  [width=broad, frame=off, background=color, backgroundcolor=panelbg,
%   leftframe=on, framecolor=accent, rulethickness=2pt, offset=1.2em]
%
%\defineframedtext
%  [solucion]
%  [width=broad, frame=off, background=color, backgroundcolor=panelbg,
%   leftframe=on, framecolor=accentb, rulethickness=2pt, offset=1.2em]
%
%\defineframedtext
%  [explicacion]
%  [width=broad, frame=off, background=color, backgroundcolor=panelbg,
%   leftframe=on, framecolor=accentb, rulethickness=2pt, offset=1.2em]
%
%\defineframedtext
%  [observacion]
%  [width=broad, frame=off, background=color, backgroundcolor=panelbg,
%   leftframe=on, framecolor=accent, rulethickness=2pt, offset=1.2em]
%
%\defineframedtext
%  [resultado]
%  [width=broad, frame=off, background=color, backgroundcolor=resultbg,
%   leftframe=on, framecolor=resultgreen, rulethickness=2pt, offset=1.2em]
%
%\define[1]\Problema{\blank[small]\startproblema{\bf Problema. }#1\stopproblema\blank[small]}
%\define[1]\Solucion{\blank[small]\startsolucion{\bf Solución. }#1\stopsolucion\blank[small]}
%\define[1]\Explicacion{\blank[small]\startexplicacion{\bf Explicación. }#1\stopexplicacion\blank[small]}
%\define[1]\Observacion{\blank[small]\startobservacion{\bf Observación. }#1\stopobservacion\blank[small]}
%\define[1]\Resultado{\blank[small]\startresultado{\bf\color[resultgreen]{Resultado.} }#1\stopresultado\blank[small]}
%
%% ---------------------------------------------------------------------
%% 6. Interacción
%% ---------------------------------------------------------------------
%\setupinteraction [state=start, color=accentb, contrastcolor=accentb,
%                   title={Tablas de contenido en ConTeXt}]
%
%% ---------------------------------------------------------------------
%% 7. Pie de página
%% ---------------------------------------------------------------------
%\setupheadertexts  [][]
%\setupfootertexts  [][{\ss\small\color[softgray]{\pagenumber}}]

\environment env_tuto
% =====================================================================
\starttext

% ------------------------ PORTADA -------------------------------------
\startstandardmakeup[pagestate=start]
  \blank[3*big]
  \startcolor[tutorialaccent]\hrule height 2pt\stopcolor
  \blank[medium]
  {\ss\bf\switchtobodyfont[24pt]\color[inkblue]{Tutorial}}\par
  \blank[small]
  {\ss\bf\switchtobodyfont[17pt]\color[inkblue]{Domina las tablas de contenido}}\par
  {\ss\bf\switchtobodyfont[17pt]\color[tutorialaccent]{avanzadas en ConTeXt}}\par
  \blank[medium]
  \startcolor[accentb]\hrule height 0.8pt\stopcolor
  \blank[2*big]
  {\it\color[softgray]{TdM completas o parciales: formato, sangría,
   entradas manuales y secciones no numeradas.}}
  \vfill
  {\ss\small\color[softgray]{Guía práctica paso a paso}}
\stopstandardmakeup

% ------------------------ RESUMEN ---------------------------------------
\startfrontmatter

\tutorialsection{Resumen}
En este tutorial aprenderemos a crear y personalizar tablas de
contenido avanzadas en ConTeXt. Al final, podrás generar TdM
completas o parciales, controlar su formato y añadir entradas
manualmente.

\blank[medium]
{\ss\bf\color[inkblue]{Contenido}}
\startcolor[codeframe]\hrule\stopcolor
\blank[small]
\placelist[tutorialsection][criterium=all,alternative=c]

\stopfrontmatter

% ------------------------ CUERPO ----------------------------------------
\startbodymatter

% =====================================================================
\tutorialsection{Prepara tu espacio de trabajo}

Primero, crea una carpeta llamada \tt{mi\_proyecto\_toc\_avanzado} y
dentro de ella, un archivo llamado \tt{documento.tex} con este
contenido inicial:

\blank[small]
\startcontextcod
\starttext

\chapter{Primer capítulo}
\section{Primera sección}
\subsection{Subsección 1.1}

\chapter{Segundo capítulo}
\section{Segunda sección}
\subsection{Subsección 2.1}

\stoptext
\stopcontextcod
\blank[small]

Compila con \tt{context documento.tex} y observa el resultado.
Obtendrás un documento con capítulos y secciones. Ahora, vamos a
añadir una tabla de contenido avanzada.

% =====================================================================
\tutorialsection{Crea una tabla de contenido con formato personalizado}

\Problema{Generar una TdM con números de página alineados a la derecha
y puntos de separación.}

\Solucion{Usa \type{\completecontent} con la alternativa \type{c}:}

\blank[small]
\startcontextcod
\starttext

\completecontent[alternative=c]

\chapter{Primer capítulo}
\section{Primera sección}
\subsection{Subsección 1.1}

\chapter{Segundo capítulo}
\section{Segunda sección}
\subsection{Subsección 2.1}

\stoptext
\stopcontextcod
\blank[small]

\Resultado{La TdM ahora muestra los números de página alineados a la
derecha con puntos de separación.}

\Explicacion{La alternativa \type{c} crea una TdM con puntos entre el
título y el número de página.}

% =====================================================================
\tutorialsection{Controla la profundidad de la tabla de contenido}

\Problema{Mostrar solo capítulos y secciones, sin subsecciones.}

\Solucion{Usa \type{\setupcombinedlist} para limitar la profundidad:}

\blank[small]
\startcontextcod
\setupcombinedlist[content][list={chapter,section}]

\starttext

\completecontent

\chapter{Primer capítulo}
\section{Primera sección}
\subsection{Subsección 1.1}

\chapter{Segundo capítulo}
\section{Segunda sección}
\subsection{Subsección 2.1}

\stoptext
\stopcontextcod
\blank[small]

\Resultado{La TdM muestra capítulos y secciones, pero no
subsecciones.}

\Observacion{\type{list=\{chapter,section\}} especifica qué niveles
de sección incluir.}

% =====================================================================
\tutorialsection{Personaliza el estilo de una entrada de la TdM}

\Problema{Cambiar el color y estilo de los títulos de sección en la
TdM.}

\Solucion{Usa \type{\setuplist} para configurar las entradas:}

\blank[small]
\startcontextcod
\setuplist[section][style={\it}, color=darkblue]

\starttext

\completecontent[alternative=c]

\chapter{Primer capítulo}
\section{Primera sección}
\subsection{Subsección 1.1}

\chapter{Segundo capítulo}
\section{Segunda sección}
\subsection{Subsección 2.1}

\stoptext
\stopcontextcod
\blank[small]

\Resultado{Los títulos de sección aparecen en cursiva y azul oscuro
en la TdM.}

\Explicacion{\type{\setuplist} permite formatear entradas de la TdM
de forma independiente.}

% =====================================================================
\tutorialsection{Ajusta la sangría de las entradas}

\Problema{Aumentar la sangría de las subsecciones en la TdM.}

\Solucion{Usa la opción \type{margin} en \type{\setuplist}:}

\blank[small]
\startcontextcod
\setuplist[subsection][margin=2em]

\starttext

\completecontent[alternative=c]

\chapter{Primer capítulo}
\section{Primera sección}
\subsection{Subsección 1.1}

\chapter{Segundo capítulo}
\section{Segunda sección}
\subsection{Subsección 2.1}

\stoptext
\stopcontextcod
\blank[small]

\Resultado{Las subsecciones aparecen con mayor sangría en la TdM.}

% =====================================================================
\tutorialsection{Incluye secciones no numeradas en la TdM}

\Problema{Que una sección no numerada (como un prólogo) aparezca en
la TdM.}

\Solucion{Usa \type{\setuphead} para activar \type{incrementnumber}
y \type{\setupcombinedlist} para incluirla:}

\blank[small]
\startcontextcod
\setuphead[subject][incrementnumber=yes]
\setupcombinedlist[content][list={chapter,section,subject}]

\starttext

\completecontent

\chapter{Primer capítulo}
\section{Primera sección}
\startsubject[title=Prólogo no numerado]
\stopsubject

\stoptext
\stopcontextcod
\blank[small]

\Resultado{El prólogo aparece en la TdM, aunque no esté numerado en
el cuerpo.}

% =====================================================================
\tutorialsection{Añade una entrada manual a la TdM}

\Problema{Insertar una entrada en la TdM que no corresponde a una
sección real.}

\Solucion{Usa \type{\writetolist} en cualquier punto del documento:}

\blank[small]
\startcontextcod
\starttext

\completecontent

\chapter{Primer capítulo}
\section{Primera sección}

\writetolist[section]{}{Entrada manual}

\section{Segunda sección}

\stoptext
\stopcontextcod
\blank[small]

\Resultado{Aparece una entrada ``Entrada manual'' en la TdM entre las
secciones.}

\Observacion{El segundo y tercer argumento de \type{\writetolist} no
siempre funcionan como se espera, pero el texto sí se añade.}

% =====================================================================
\tutorialsection{Crea una tabla de contenido parcial en un capítulo}

\Problema{Mostrar solo las secciones de un capítulo al inicio del
mismo.}

\Solucion{Coloca \type{\placecontent} dentro del capítulo:}

\blank[small]
\startcontextcod
\starttext

\chapter{Primer capítulo}
\placecontent

\section{Primera sección}
\subsection{Subsección 1.1}
\section{Segunda sección}

\chapter{Segundo capítulo}
\section{Tercera sección}

\stoptext
\stopcontextcod
\blank[small]

\Resultado{Dentro del primer capítulo aparece una TdM con solo sus
secciones.}

% =====================================================================
\tutorialsection{Personaliza el título de la TdM}

\Problema{Cambiar ``Tabla de contenido'' por ``Índice general''.}

\Solucion{Usa \type{\setupheadtext} para modificar el texto:}

\blank[small]
\startcontextcod
\setupheadtext[es][content={Índice general}]

\starttext

\completecontent

\chapter{Primer capítulo}
\section{Primera sección}

\stoptext
\stopcontextcod
\blank[small]

\Resultado{El título de la TdM ahora es ``Índice general''.}

% =====================================================================
\tutorialsection{Prepara un documento final con TdM avanzada}

\Problema{Usar todas las técnicas avanzadas en un documento
coherente.}

\Solucion{Crea un documento que combine formato personalizado,
sangría, entradas manuales y secciones no numeradas:}

\blank[small]
\startcontextcod
% Configuraciones globales
\setupheadtext[es][content={Índice general}]
\setupcombinedlist[content][list={chapter,section,subject}]

% Formato de entradas
\setuplist[section][style={\it}, color=darkblue]
\setuplist[subsection][margin=2em, color=darkgreen]
\setuplist[subject][color=darkred]

% Incluir sujetos no numerados
\setuphead[subject][incrementnumber=yes]

\starttext

\completecontent[alternative=c]

\chapter{Introducción}
\section{Primera sección}

\chapter{Desarrollo}
\placecontent  % TdM parcial del capítulo

\section{Segunda sección}
\subsection{Subsección 2.1}

\startsubject[title=Tema especial]
\stopsubject

\section{Tercera sección}

\writetolist[section]{}{Nota adicional}

\chapter{Conclusiones}
\section{Cuarta sección}

\stoptext
\stopcontextcod
\blank[small]

\Resultado{El documento tiene una TdM global con formato
personalizado, secciones no numeradas, una entrada manual y una TdM
parcial en el capítulo ``Desarrollo''.}
\tutorialsection{Ejercicio}

Practica lo aprendido intentando una tabla de contenidos como la siguiente:
% ============================================================
% DOCUMENTO DE MUESTRA
% Índice de contenido — versión estable
% ConTeXt LMTX
% ============================================================

\setuppapersize[letter]
\mainlanguage[es]

% ------------------------------------------------------------
% TIPOGRAFÍA
% ------------------------------------------------------------

\definefontfamily[mainface][rm][Pagella]
\definefontfamily[mainface][ss][Iwona]
\definefontfamily[mainface][tt][ModernMono]

\setupbodyfont[mainface,11pt]

% ------------------------------------------------------------
% COLORES
% ------------------------------------------------------------

\definecolor[Azul][h=315B7D]
\definecolor[AzulOscuro][h=173A5A]
\definecolor[Rojo][h=B33A32]
\definecolor[Gris][h=687078]
\definecolor[Fondo][h=F4F6F7]

% ------------------------------------------------------------
% PÁGINA
% ------------------------------------------------------------

\setuplayout
  [
    width=middle,
    height=middle,
    backspace=24mm,
    topspace=18mm,
    header=8mm,
    footer=8mm,
    headerdistance=4mm,
    footerdistance=4mm
  ]

% ------------------------------------------------------------
% CAPÍTULOS Y SECCIONES
% ------------------------------------------------------------

\setuphead
  [chapter,section,subsection]
  [page=no]

\setupcombinedlist
  [content]
  [list={chapter,section}]

% ------------------------------------------------------------
% ENCABEZADO Y PIE
% ------------------------------------------------------------

%\setupheadertexts
%  []
%  [{\switchtobodyfont[8pt]\ss\color[Gris]
%    CUADERNO DE MÉTODOS\, \diamond\ MATEMÁTICAS}]

%\setupfootertexts
%  [{\switchtobodyfont[8pt]\ss\color[Gris] MALAQUÍAS}]
%  [{\switchtobodyfont[8pt]\ss\color[Gris] \pagenumber}]

% ------------------------------------------------------------
% CONFIGURACIÓN DEL ÍNDICE
% ------------------------------------------------------------

\setuplist
  [chapter]
  [
    alternative=b,
    width=10mm,
    distance=1mm,
    textstyle={\switchtobodyfont[14pt]\ss\bf},
    textcolor=AzulOscuro,
    pagenumberstyle={\switchtobodyfont[10pt]\ss},
    pagenumbercolor=Gris
  ]

\setuplist
  [section]
  [
    alternative=b,
    width=10mm,
    distance=1mm,
    textstyle={\switchtobodyfont[11pt]\ss},
    textcolor=Gris,
    pagenumberstyle={\switchtobodyfont[9pt]\ss},
    pagenumbercolor=Gris
  ]

\setuplist
  [subsection]
  [
    alternative=b,
    width=10mm,
    distance=1mm,
    textstyle={\switchtobodyfont[9pt]\ss\it},
    textcolor=Gris,
    pagenumberstyle={\switchtobodyfont[8pt]\ss},
    pagenumbercolor=Gris
  ]

% ------------------------------------------------------------
% RÓTULO GENÉRICO
% ------------------------------------------------------------

\define[2]\Rotulo
  {{\switchtobodyfont[#1]#2}}

% ============================================================
% DOCUMENTO
% ============================================================


% ============================================================
% PORTADA DEL ÍNDICE
% ============================================================

\Rotulo{10pt}{CUADERNO DE MÉTODOS}

\blank[small]

\Rotulo{24pt}{Mapa del libro}

\blank[small]

\Rotulo{11pt}{Una arquitectura de ideas, métodos y problemas}

\blank[medium]

\blackrule
  [
    width=\textwidth,
    height=1.2pt,
    color=AzulOscuro
  ]

\blank[big]

\Rotulo{11pt}{
  Este libro está organizado como una progresión:
  primero se reconoce la estructura del problema;
  después se elige el método; finalmente se ejecuta
  el cálculo.
}

\blank[big]

\blackrule
  [
    width=30mm,
    height=1pt,
    color=Rojo,
  ]

\blank[small]
\completecontent[criterium=all]
% ============================================================
% CONTENIDO DE MUESTRA
% ============================================================

\startchapter[title={El método de intervalos}]

  \startsection[title={Signos y puntos críticos}]

    \startsubsection[title={Definición de punto crítico}]

      Un punto crítico es aquel donde la derivada se anula
      o no existe.

    \stopsubsection

    \startsubsection[title={Signos en cada intervalo}]

      Se evalúa la función en un punto de prueba dentro
      de cada intervalo.

    \stopsubsection

  \stopsection

  \startsection[title={Construcción de la tabla de signos}]

    La tabla organiza los intervalos y los signos resultantes.

  \stopsection

\stopchapter

% ============================================================
% ÍNDICE
% ============================================================

%\completecontent

% ============================================================
% CIERRE EDITORIAL
% ============================================================

\blank[big]

\blackrule
  [
    width=30mm,
    height=1pt,
    color=Rojo,
  ]

\blank[small]

%\Rotulo{8pt}{MÉTODOS \texttbullet ESTRUCTURA \textbullet RESOLUCIÓN}

\stopbodymatter

% ------------------------ APÉNDICES --------------------------------------
\startappendices

\tutorialsection{¡Has logrado!}

Dominar las tablas de contenido avanzadas en ConTeXt. Ahora sabes:

\startitemize
\item \bf{Controlar el formato}: usar \type{alternative} para
      diferentes estilos de TdM.
\item \bf{Limitar la profundidad}: usar \type{\setupcombinedlist}
      para seleccionar niveles.
\item \bf{Personalizar entradas}: usar \type{\setuplist} para estilo,
      color y sangría.
\item \bf{Incluir secciones no numeradas}: usar \type{\setuphead} y
      \type{\setupcombinedlist}.
\item \bf{Añadir entradas manuales}: usar \type{\writetolist} para
      insertar elementos.
\item \bf{Crear TdM parciales}: usar \type{\placecontent} dentro de
      capítulos.
\item \bf{Personalizar el título}: usar \type{\setupheadtext} para
      cambiar el texto.
\stopitemize

Este flujo de trabajo te permite crear tablas de contenido adaptadas
a las necesidades de tus documentos.

\stopappendices

% ------------------------ BACKMATTER --------------------------------------
\startbackmatter

\tutorialsection{Próximos pasos}

\startitemize
\item \bf{Explora más opciones}: \type{\setuplist} tiene opciones
      como \type{numberalign}, \type{width}, \type{distance}.
\item \bf{Combina con otros elementos}: usa TdM en diferentes
      idiomas o con diferentes estilos.
\item \bf{Crea plantillas}: define configuraciones de TdM en un
      archivo separado.
\item \bf{Recuerda}: una buena TdM mejora la navegación y comprensión
      del documento.
\stopitemize

\stopbackmatter

\stoptext